\scanpage{192}
\mspage{66}
% archive sheet number: 191

\uline{Définition} 5.12. 1) Soit $\mathcal M$ une catégorie satisfaisant les conditions a) b) c) d) (p. 44.\unsure{--|$=$}45), \inserted{soit $\unsure{\pi|\Pi}\geq\pi_0$ un card\unsure{inal}.} Une partie $M$ de $\mathrm{Ob}\,\mathcal M$ est dite \uline{$L$-$\pi$-accessible}, si elle est $\pi$-accessible et stable par \inserted{\uline{petites}} $\underline{\lim}$ filtrantes\unsure{*}, ou ce qui revient au $=$ \inserted{(5.9)}, si elle est $\pi$-accessible et si $M^\pi=M\cap\mathrm{Filt}^{\Pi}$ est stable par $\underline{\lim}$ filtrantes de cardinal $\leq\pi$.

% slanted
\margin{On dit alors que $M$ est \uline{$L$-stable}}

(\uline{NB} Si les objets de $\mathrm{Filt}^{\Pi}$ sont $\pi$-accessibles, 5.8 nous dit qu'il cela signifie aussi qu'il \inserted{suffit pour cela qu'il} existe une partie $M_0$ de $\mathrm{Filt}^{\Pi}$ telle que $M$ soit la \inserted{$L$-}enveloppe (enveloppe stable par $\underline{\lim}$ filtrantes) de $M_0$ -- condition qui est de \unsure{ttes|tte} \unsure{façons|façon} tjrs nécessaire, en prenant $M_0=M^\pi$.)

2) Soit $M$ une partie de $\mathrm{Ob}\,\mathcal M$. On \deleted{\unsure{dit elle}} qu'elle est \uline{$L$-accessible}, si elle satisfait à l'une des \deleted{deux} \inserted{trois} conditions équivalentes suivantes :

a) Il existe un cardinal $\pi$ tel que $M$ soit $L$-$\pi$-accessible, \deleted{\unsure{b).}}

\inserted{a$'$) $M$ est accessible et $L$-stable.}

b) Il existe une petite partie $M_0$ de $\mathrm{Ob}\,\mathcal M$, telle que $M$ soit la $L$-enveloppe de $M_0$.

% slanted
\margin{$\longleftarrow$ Cette condition ne dépend pas de la filtration cardinale choisie de $\mathcal M$}

b$'$) Il existe une partie $M_0$ de $\mathcal M$ contenue dans une partie $\mathrm{Filt}^{\pi}$ (ou ce qui revient au $=$, telle que la \inserted{sous} catégorie pleine \deleted{\unsure{de}} $\underline M_0$ de $\mathcal M$ soit \unsure{ess\textsuperscript{t.}} petite, i.e. égale à une petite catégorie), telle que $M$ soit la $L$-enveloppe de $M_0$.

